\documentclass{article}
\usepackage[T1]{fontenc}
\usepackage{lmodern}
\usepackage{graphicx}
\usepackage{amsmath,amssymb}
\usepackage[a4paper,margin=2cm]{geometry}
\usepackage[absolute,overlay]{textpos}
\setlength{\TPHorizModule}{1mm}
\setlength{\TPVertModule}{1mm}
\usepackage[indonesian]{babel}
\usepackage{hyperref}
\setlength{\emergencystretch}{2em}
% unit: Halaman 64

\begin{document}

% === Halaman 64 (llm) ===

\textcolor{red}{\textbf{2. Auto-Key Vig\`en\`ere cipher}}

\begin{itemize}
    \item Jika panjang kunci lebih kecil dari panjang plainteks, maka kunci disambung dengan plainteks tersebut.
    \item Misalnya,
    \begin{itemize}
        \item Pesan: \texttt{negara penghasil minyak mentah di dunia}
        \item Kunci: \texttt{INDO}
    \end{itemize}
    maka kunci tersebut disambung dengan plainteks semula sehingga panjang kunci menjadi sama dengan panjang plainteks:
    \begin{itemize}
        \item Plainteks: \texttt{negarapenghasilminyakmentahdidunia}
        \item Kunci: \texttt{INDONEGARAPENGHASILMINYAKMENTAHDID}
        \item Cipherteks: \texttt{VRJOEEVEEGWEFOSMAVJMSZCNDMLQBDBQQD}
    \end{itemize}
\end{itemize}

\end{document}
